%% kapa: exam=1 round=28 year=2017 period=2 form=A subject=appraisal_law subject_ko=감정평가관계법규 seq=26 number=26
%% source: Q-Net 감정평가사 기출문제 (cst00309) artlSeq=5199603 file=extracted/1008/5199603/2017년 제28회 감정평가사 1차시험 문제지/2교시 A형.pdf
%% conversion: kordoc 4.24.51 -> tools/gen1.py (auto); review=proofread-sonnet; flags=none
\begin{question}{26}
건축법령상 공개공지등의 확보에 관한 설명으로 옳지 않은 것은?
\begin{choices}
  \item 상업지역에서 업무시설로서 해당 용도로 쓰는 바닥면적의 합계가 5천 제곱미터 이상인 건축물의 대지에는 공개공지등을 확보하여야 한다.
  \item 공개 공지는 필로티의 구조로 설치할 수 있다.
  \item 공개공지등에는 물건을 쌓아 놓거나 출입을 차단하는 시설을 설치하지 아니하여야 한다.
  \item 공개공지등의 면적은 건축면적의 100분의 10 이하로 한다.
  \item 공개공지등을 설치하는 경우에는 건축물의 용적률 기준을 완화하여 적용할 수 있다.
\end{choices}
\end{question}
\answer{4}
